\subsection{Canonical two-adic norms and the simultaneous radical fields}

\begin{statement}{A.101}{Proved}
{The actual permutation Kummer field and its boundary intersections}
Let $a=3u^8$, $b=-2v^8$ be the current coefficients of
A.44, with $u,v$ prime to six, and let $N=3^n$, $n\ge1$.
Put $K=\Q(\alpha)$, $\alpha^N=a$.
Let $\beta\in K^\times$ have rational norm $b$ and the
following specified divisor above two: valuation one
at the unique degree-one prime, and zero at every
other prime of $K$ above two.
Fix $r\ge1$, an integer $M$ divisible by $\operatorname{lcm}(24,N)$, and
a cyclotomic field
\[
 R=\Q(\mu_D),\qquad \operatorname{lcm}(M,3^r)\mid D.
 \tag{A.101.1}
\]
Write $L_R=R(\alpha)$ and let $\tau(\alpha)=\zeta_N\alpha$,
$\tau|_R=1$.
For $i\in\Z/N\Z$ put $\beta_i=\tau^i\beta$, and choose
$\delta_i^{3^r}=\beta_i$.
Define the actual splitting fields
\[
 S_\beta=L_R(\delta_i:i\in\Z/N\Z),\qquad
 B_R=R(b^{1/M}),\qquad S=S_\beta B_R.
 \tag{A.101.2}
\]
Then
\[
 \operatorname{Gal}(S_\beta/R)
          =(\Z/3^r\Z)^N\rtimes\Z/N\Z
 \tag{A.101.3}
\]
with the literal cyclic permutation action, and with
the explicit splitting on the displayed generators
given below.
Its maximal abelian subfield over $R$ is
\[
 (S_\beta)^{\rm ab/R}
      =R(\alpha,\delta_*),\qquad
      \delta_*=\prod_i\delta_i,\qquad \delta_*^{3^r}=b.
 \tag{A.101.4}
\]
Put $m=v_3(M)$ and $d=\min(r,m)$.
The exact boundary intersection is
\[
                  S_\beta\cap B_R=R(b^{1/3^d}).
 \tag{A.101.5}
\]
Both $S_\beta$ and the simultaneous field $S$ have
maximal abelian subfield over $\Q$ exactly $R$.
The assertion for $S$ follows from the full simultaneous
group computation below, including its common
$b$-subfield.

Let $F$ be the compositum of
$L=\Q(\alpha,\mu_N)$ with the conjugates of a three-primary
ray class field of $K$ whose modulus excludes two,
as in A.99.  Then
\[
                         FR\cap S=L_R.
 \tag{A.101.6}
\]
Consequently the permitted split primes generating
the class targets in A.95 and A.100 may additionally
split every conjugate radical in $S_\beta$.
Any larger cyclotomic CRT automorphism fixing $R$
can be retained.
In particular, if $v_3(D)=r\ge v_3(M)$, the additional
condition $v_3(q-1)=r$ is available by prescribing the
nontrivial class $1+3^r$ modulo $3^{r+1}$.
\end{statement}
\begin{proof}
\emph{The valuation matrix and the literal automorphisms.}
The $N$th power map on $\Z_2^\times$ is an automorphism:
it is bijective on its two-primary torsion, and Hensel
lifting gives the unique root in each unit residue
class for the odd exponent $N$.
Thus $a$ has a unique $N$th root in $\Q_2$.
All its other roots lie in $\Q_2(\mu_N)$.
The root algebra over $\Q_2$ is unramified, since
$a$ and $N$ are units at two; it has exactly one
degree-one factor.

Since $R$ contains $\mu_N$, the extension $L_R/R$
splits completely at every prime above two.
Its degree is $N$: the full radical degree over any
abelian cyclotomic extension was proved in A.44/A.77
using $v_3(a)=1$.
Fix a prime $w\mid2$ of $R$, and label its $N$ primes
in $L_R$ compatibly with the conjugates $\beta_i$.
If $e_w=e(R_w/\Q_2)$, the assumed valuation pattern
gives the exact matrix
\[
                       v_{w_j}(\beta_i)=e_w\delta_{ij}.
 \tag{A.101.7}
\]
Here $e_w$ is a power of two.
The entries are therefore three-adic units on the
diagonal.  For every $h\ge1$, the valuations make the
classes of the $\beta_i$ independent in
$L_R^\times/(L_R^\times)^{3^h}$.
Indeed a power relation forces each exponent to be
zero modulo $3^h$ by (A.101.7).
Kummer theory now gives the full degree $3^{rN}$
over $L_R$ and the independent automorphisms
$\delta_i\mapsto\zeta_{3^r}^{z_i}\delta_i$.

This degree statement also justifies actual lifts,
rather than only an abstract action on a module.
The field is the tensor root algebra with the independent
relations $\delta_i^{3^r}=\beta_i$ over $L_R$.
The assignments
\[
 \widetilde\tau(\alpha)=\zeta_N\alpha,\quad
 \widetilde\tau(\delta_i)=\delta_{i+1},\quad
 \widetilde\tau|_R=1
 \tag{A.101.8}
\]
preserve those relations and define an automorphism.
It has order exactly $N$ and proves the splitting in
(A.101.3).
For $u\in\operatorname{Gal}(R/\Q)$, use its cyclotomic
exponent modulo $D$ and define
\[
 \widetilde u|_R=u,\quad \widetilde u(\alpha)=\alpha,
                \quad\widetilde u(\delta_i)=\delta_{ui}.
 \tag{A.101.9}
\]
These maps again preserve every relation.
They satisfy $\widetilde u\widetilde v=\widetilde{uv}$
and $\widetilde u\widetilde\tau\widetilde u^{-1}
=\widetilde\tau^u$ on all the generators.
On the Kummer kernel their conjugation action is
the permutation of the indices together with
multiplication by the cyclotomic exponent $u$
on $\Z/3^r\Z$.
This proves the full arithmetic action explicitly.

Write $V=(\Z/3^r\Z)^N$ and $T=\Z/N\Z$.
The image of $\widetilde\tau-1$ on the permutation
module $V$ is precisely the kernel of the sum map.
Thus the commutator subgroup of $V\rtimes T$ is
$\ker(\sum:V\to\Z/3^r\Z)$, and its abelianization
is $T\oplus\Z/3^r\Z$.
The fields representing these two quotient maps are
$R(\alpha)$ and $R(\delta_*)$, respectively:
$\widetilde\tau$ fixes $\delta_*$ and a vector $(z_i)$
sends it to $\zeta_{3^r}^{\sum z_i}\delta_*$.
The equality $\prod_i\beta_i=b$ is the actual field
norm of $\beta$, so this proves (A.101.4).

Over $\Q$, the $T$-coinvariants of $V$ are one copy
of $\Z/3^r\Z$ with cyclotomic action.
The automorphism in (A.101.9) with $u=-1$ acts there
by $-1$, so its coinvariants are zero.
The same $u$ acts by $-1$ on $T$, whose coinvariants
are also zero.
The entire kernel over $R$ is therefore killed by
the arithmetic abelianization.
Since $R/\Q$ is already abelian, the maximal abelian
field of $S_\beta/\Q$ is exactly $R$.

\emph{The boundary intersection and its simultaneous group.}
The extension $B_R/R$ is cyclic and its three-primary
subextension is $R(b^{1/3^m})$.
This subextension has degree $3^m$ and is totally
ramified at every prime above two, by $v_2(b)=1$
and the coprime valuation calculation (A.95.6).
The remaining degree is prime to three.
Because $S_\beta/R$ is a three-extension, its intersection
with $B_R$ is in this three-primary part.
It is an abelian extension of $R$, so (A.101.4)
places it in $R(\alpha,b^{1/3^r})$.
If $m>r$, the extension obtained by adjoining $\alpha$
over $R(b^{1/3^r})$ is unramified at two, whereas
$R(b^{1/3^m})/R(b^{1/3^r})$ is totally ramified.
Their intersection is therefore the base field.
For $m\le r$ the indicated $b$-root is already in
$S_\beta$.  This proves (A.101.5).

We give the simultaneous arithmetic computation.
Let $S_3=S_\beta R(b^{1/3^m})$ and choose
$\eta^{3^m}=b$ with
\[
              \delta_*^{3^{r-d}}=\eta^{3^{m-d}}.
 \tag{A.101.10}
\]
Such a choice is obtained by multiplying a root by
a root of unity in $R$.
Equation (A.101.5) proves that (A.101.10) is exactly
the common Kummer relation.
The actual group over $L_R$ is therefore
\[
 V'=\left\{(z_0,\ldots,z_{N-1},w)
        \in(\Z/3^r\Z)^N\oplus\Z/3^m\Z:
            \sum_i z_i\equiv w\pmod{3^d}\right\}.
 \tag{A.101.11}
\]
The lifts (A.101.8)--(A.101.9) extend by fixing
$\eta$; they preserve (A.101.10), so these are
literal automorphisms of the simultaneous field.
The translation lift cyclically permutes the $z_i$
and fixes $w$.
Its coinvariants on $V'$ are exactly
\[
 \{(z,w)\in\Z/3^r\Z\oplus\Z/3^m\Z:
                           z\equiv w\pmod{3^d}\}
                       \simeq\Z/3^{\max(r,m)}\Z.
 \tag{A.101.12}
\]
To verify this map, use $(z_i,w)\mapsto(\sum z_i,w)$.
The projection $V'\to(\Z/3^r\Z)^N$ is onto, and
the differences under cyclic permutation are exactly
the sum-zero vectors with $w=0$.
The two coordinates in (A.101.12) are compatible
reductions of the one with larger modulus.
Cyclotomic $u=-1$ acts by $-1$ on this quotient,
and also on the translation group $T$.
This proves that the maximal $\Q$-abelian field
of $S_3$ is $R$.

Let $B_{3'}$ be the prime-to-three subextension of
$B_R/R$.
It is Galois over $\Q$, since its subgroup over $R$
is characteristic in that cyclic group.
Its intersection with $S_3$ is $R$, so the group
over $R$ of $S=S_3B_{3'}$ is the direct product
of its three-primary group and this prime-to-three
cyclic group.
The conjugation actions on the three-primary group
are exactly the actions computed above, hence its
image in every arithmetic abelian quotient is zero.
The remaining quotient factors through
$\operatorname{Gal}(B_{3'}/\Q)$.
This last field has no abelian subfield beyond $R$.
Indeed, A.44 gives
$\Q(\mu_M,b^{1/M})^{\rm ab}=\Q(\mu_M)$.
For its compositum with $R$, the arithmetic group
is the fibre product over $\Q(\mu_M)$.
Its commutator subgroup is exactly the pairs
$(g',1)$ with $g'$ in the original commutator
subgroup: the second factor is abelian and projection
onto the first group is onto.
Its maximal abelian field is therefore exactly $R$.
The same holds for the subfield $B_{3'}$ containing
$R$.  This proves the asserted maximal abelian
field of the full simultaneous $S$ without inferring
it from the two individual abelianizations.

\emph{The ray-Hilbert intersection.}
The field $FR/R$ is a three-extension unramified at
two, and it contains $L_R$.
After adjoining the prime-to-three field $B_{3'}$
to the base $L_R$, it is still an unramified
three-extension at all primes above two.
This base change has ramification index prime to
three there.
The matrix (A.101.7) therefore remains diagonal
with three-adic unit entries.

Put $h=\max(r,m)$.
Over this base the Kummer subgroup for the simultaneous
three-extension is generated by the radicands
\[
          \beta_i^{3^{h-r}}\quad(0\le i<N),
                   \qquad b^{3^{h-m}}
                          =\prod_i\beta_i^{3^{h-m}}
 \tag{A.101.13}
\]
modulo $3^h$th powers.
The valuation matrix embeds this subgroup into
$(\Z/3^h\Z)^N$.
In particular every nontrivial cyclic Kummer
subextension is ramified at some prime above two:
a nonzero vector of order $3^j$ has a coordinate
of that exact order, and the corresponding root
has a valuation denominator $3^j$ there.
Thus this simultaneous abelian three-extension
has no nontrivial subextension unramified at
all the primes above two.
It has trivial intersection with the base-changed
ray-Hilbert extension.
The prime-to-three base change itself has trivial
intersection with $FR/L_R$ by the coprime degrees.
Descending these two intersection statements gives
exactly (A.101.6).

Finally let $R'$ be a larger cyclotomic field and
let its CRT automorphism fix $R$.
The proved maximal abelian field gives $S\cap R'=R$,
so the identity on $S$ and that automorphism on $R'$
are compatible.
The proof of (A.101.6) works after replacing $R$
by $R'$, since its two-adic ramification index
is still prime to three.
Thus the compatible Hilbert Frobenius coset of
A.95 combines with the identity on $S$.
Its exact coset-character argument preserves the
same split-prime class generators; with the subgroup
$J$ it also preserves the relative generators
in A.100.
Chebotarev supplies all these primes outside any
finite exceptional set.
Taking the cyclotomic class $1+3^r$ modulo
$3^{r+1}$ gives the stated exact valuation of $q-1$.
\end{proof}

The canonical valuation hypothesis is available by
finite weak approximation on the norm fibre, using
the norm-one torus in A.77 and
\cite[Theorems 1.2, 2.2 and 3.2]{RootHoshiYamasaki2025}.
At two, the degree-one factor can carry $b$ and
the other factors can carry one, giving a local
norm element with exactly the required divisor.
The repairs in A.99 and A.100 retain this divisor
because two is included in their controlled set.
For an existing exact finite norm chain ending at level $t$,
apply the theorem with $n=t$, with $\beta$ equal to its
top representative, and with
$L=\Q(a^{1/3^t},\mu_{3^t})$.
The conjugates of every lower representative are block
products of the top conjugates under the actual
field norm.  Their equal-precision radicals are
therefore already in the displayed top radical
field.  These facts control the auxiliary prime
conditions; they do not yet solve the new local
unit norm equations at primes where a repaired
representative has acquired nonzero valuations.

\noindent The exact certificate
\path{code/norm_canonical_wreath_certificate.py}, with retained output
\path{code/norm_canonical_wreath_certificate_output.txt}, verifies
the arithmetic permutation action with its cyclotomic factor,
the simultaneous common-$b$ kernel and its coinvariants,
and the injective valuation lattices used in the intersection
proof.  The actual field automorphisms are defined in
(A.101.8)--(A.101.10), using the proved Kummer degree.
